\chapter{Laplace Equation}

\section{Mean value property}

\begin{definition}
A $C^2$ function satisfying $\Delta u=0$ is called \emph{harmonic}.
\end{definition}

\begin{theorem}[Mean value property]
If $u\in C^2(U)$ is harmonic, then
$$
u(x)= \fint_{\partial B(x,r)} u(x)\,dS
= \fint_{B(x,r)} u(x)\,dx
$$
for each $B(x,r)\subset U$.
\end{theorem}

\begin{theorem}[Converse to the mean value property]
If $u\in C^2(U)$ satisfies
$$
u(x)= \fint_{\partial B(x,r)} u(x)\,dS
= \fint_{B(x,r)} u(x)\,dx
$$
for $B(x,r)\subset U$,
then $u$ is harmonic.
\end{theorem}


\section{Maximum principle, positivity, uniqueness, Harnack inequality}

\begin{definition}
    $
    C(\overline{U})
    :=
    \left\{
    u\in C(U): u \text{ is uniformly continuous on every bounded subset of } U
    \right\}.
    $
    This mmeans $u$ extends continuously to $\overline{U}$.
    \end{definition}

\begin{theorem}[Maximum principle]
Suppose $u\in C^2(U)\cap C(\overline{U})$ is harmonic within $U$.

\begin{enumerate}
    \item Then
    $$
    \max_{\overline{U}} u = \max_{\partial U} u.
    $$
    This is the \emph{weak maximum principle}.

    \item Furthermore, let $U$ be connected. If there exists $x_0\in U$ such that
    $$
    u(x_0)=\max_{\overline{U}} u,
    $$
    then $u$ is constant within $U$.
    
    This is the \emph{strong maximum principle}.
\end{enumerate}
\end{theorem}

\begin{theorem}[Positivity]
Let $U$ be open, bounded, and connected, with $U\subset \mathbb{R}^n$, and let
$
u\in C^2(U)\cap C(\overline{U}).
$
Suppose $u$ solves the boundary value problem
$$
\begin{cases}
\Delta u=0 & \text{in } U,\\
u=g & \text{on } \partial U,
\end{cases}
\qquad \text{with } g\ge 0.
$$
If $g(z)>0$ for some $z\in \partial U$, then
$
u(x)>0 \text{ for all } x\in U.
$
\end{theorem}

\begin{theorem}[Uniqueness]
Let $g\in C(\partial U)$ and $f\in C(U)$.
Then there exists at most one solution
$
u\in C^2(U)\cap C(\overline{U})
$
of
$$
\begin{cases}
-\Delta u=f & \text{in } U,\\
u=g & \text{on } \partial U.
\end{cases}
$$
\end{theorem}

\begin{definition}
$V\subset\subset U$ is defined as
$
\overline{V}\text{ is compact and }\overline{V}\subset U.
$
That is, $V$ is compactly contained in $U$.
\end{definition}

\begin{theorem}[Harnack's inequality]
For each connected open set $V\subset\subset U$, there exists a constant
$
C=C(V)>0
$
such that
$$
\sup_V u \le C \inf_V u
$$
for all nonnegative harmonic functions $u$ in $U$.
\end{theorem}

\section{Energy method and Dirichlet principle}

\begin{theorem}[Dirichlet's principle]
Assume $u\in C^2(\overline{U})$ solves
$$
\begin{cases}
-\Delta u=f & \text{in } U,\\
u=g & \text{on } \partial U,
\end{cases}
$$
the Poisson boundary value problem.
Define the energy functional
$$
I[w]:=\int_U \left( \frac12 |Dw|^2 - wf \right)\,dx.
$$
Define the admissible set
$$
\mathcal{A}:=\left\{ w\in C^2(\overline{U}) : w=g \text{ on } \partial U \right\}.
$$
Then
$$
I[u]=\min_{w\in \mathcal{A}} I[w].
$$
The converse is also true.
\end{theorem}

\section{Mathematical supplement: distributions}


\begin{definition}
For $f$,
$
\operatorname{supp}(f):=\overline{\{x\in U: f(x)\neq 0\}}.
$
\end{definition}

\begin{definition}
Define
$
\mathcal{D}(U):=
\left\{
f\in C^\infty(U): \operatorname{supp}(f)\text{ is compact}
\right\}.
$
This is the space of test functions on $U$.
$\mathcal{D}(U)$ is equipped with the following notion of convergence:
for $(f_n)\subset \mathcal{D}(U)$, $f\in \mathcal{D}(U)$,
$$
f_n\to f \text{ in } \mathcal{D}(U)
$$
means that there exists $K\subset\subset U$ such that
$
\operatorname{supp}(f_n)\subset K
$
for all $n$ and
$
\lim_{n\to\infty} \|f_n-f\|_{C^k(U)}=0
$
for all $k$, where
$$
\|f\|_{C^k(U)}
=
\sum_{|\alpha|\le k}\|D^\alpha f\|_{L^\infty(U)}.
$$
\end{definition}

\begin{claim}\
\begin{enumerate}
    \item $\mathcal{D}(U)$ is not metrizable.
    
    \item A linear functional
    $$
    T:\mathcal{D}(U)\to\mathbb{R}
    $$
    is (sequentially) continuous if and only if for every $K\subset\subset U$ there exist
    $m\in \mathbb{N}$ and $C>0$ such that, for every $\phi\in \mathcal{D}_K
    :=\{\phi\in\mathcal{D}:\operatorname{supp}(\phi)\subset K\}$,
    $$
    |T(\phi)|\le C \|\phi\|_{C^{m}(U)}.
    $$
\end{enumerate}
\end{claim}

\begin{enumerate}
    \setcounter{enumi}{2}
    \item If
    $
    f\in L^1_{\mathrm{loc}}
    :=
    \left\{
    f:U\to \mathbb{R}: f \text{ measurable},\ 
    \int_K |f|\,dx<\infty \ \forall K\subset\subset U
    \right\},
    $
    then
    $$
    T_f:\mathcal{D}(U)\to \mathbb{R},
    \qquad
    \phi\mapsto T_f(\phi)=\int_U \phi f\,dx
    $$
    is a sequentially continuous linear functional on $\mathcal{D}$, i.e. $T_f\in \mathcal{D}'$.

    \item The map
    $
    f\mapsto T_f
    $
    is injective from $L^1_{\mathrm{loc}}$ to $\mathcal{D}'$.
    So $\mathcal{D}'$ contains a copy of $L^1_{\mathrm{loc}}$.
\end{enumerate}

\begin{definition}
We call elements in $\mathcal{D}'$ \emph{distributions}.
\end{definition}

\begin{proposition}
For all multi-indices $\alpha$ and all $T\in \mathcal{D}'$,
$$
D^\alpha T:\phi\mapsto (-1)^{|\alpha|}T(D^\alpha \phi)
$$
defines a distribution $D^\alpha T\in \mathcal{D}'$.
\end{proposition}

\paragraph{Motivation.}
For $f,\phi\in C_c^\infty(U)$,
$$
\int_U (D^\alpha f)\phi\,dx
=
(-1)^{|\alpha|}
\int_U f(D^\alpha \phi)\,dx.
$$

\begin{proposition}
Suppose $u\in C^k(U)$ and $\alpha$ is a multi-index with $|\alpha|\le k$.
If
$$
T_u:\phi\mapsto \int_U u\phi\,dx
$$
is the distribution defined by $u$, then
$$
D^\alpha T_u = T_{D^\alpha u}
\qquad \text{for all } \phi\in \mathcal{D}.
$$
That is,
$$
(D^\alpha T_u)(\phi)=T_{D^\alpha u}(\phi),
$$
or equivalently,
$$
(-1)^{|\alpha|}\int_U u\,D^\alpha \phi
=
\int_U (D^\alpha u)\phi.
$$
\end{proposition}

\section{Fundamental solution and representation formulas}

\subsection*{Back to Laplace equation}

\begin{definition}
The fundamental solution of the Laplace equation is defined by
$$
\Phi(x)=
\begin{cases}
-\dfrac{1}{2\pi}\log|x|, & n=2,\\[1em]
\dfrac{1}{n(n-2)\alpha(n)}\dfrac{1}{|x|^{n-2}}, & n\ge 3.
\end{cases}
$$
\end{definition}

\begin{proposition}
It is easy to see that $\Phi(x)$ is harmonic away from $0$.
But it has exactly the right singularity at $0$ to produce a point source there:
$$
\int_{\mathbb{R}^n}\Phi(x)\Delta \phi(x)\,dx
=
-\phi(0)
\qquad \forall \phi\in C_c^2(\mathbb{R}^n).
$$
That is,
$$
\Delta T_\Phi(\phi)=-\delta_0(\phi).
$$
\end{proposition}

\begin{theorem}[Solving Poisson equation]
Define $u$ by
$$
u(x)=\int_{\mathbb{R}^n}\Phi(x-y)f(y)\,dy
=
\begin{cases}
-\dfrac{1}{2\pi}\displaystyle\int_{\mathbb{R}^2}\log|x-y|\,f(y)\,dy, & n=2,\\[1.2em]
\dfrac{1}{n(n-2)\alpha(n)}
\displaystyle\int_{\mathbb{R}^n}\dfrac{f(y)}{|x-y|^{n-2}}\,dy, & n\ge 3.
\end{cases}
$$
Then
\begin{enumerate}
    \item $u\in C^2(\mathbb{R}^n)$,
    \item $-\Delta u=f$ in $\mathbb{R}^n$.
\end{enumerate}
\end{theorem}

\begin{definition}
The Green's function for the region $U$ is
$$
G(x,y)=\Phi(y-x)-\phi^x(y),
$$
where $\phi^x=\phi^x(y)$ is called a \emph{corrector function}, which solves the boundary value problem
$$
\begin{cases}
\Delta \phi^x = 0 & \text{in } U,\\
\phi^x = \Phi(y-x) & \text{on } \partial U.
\end{cases}
$$
Here $U\subset \mathbb{R}^n$ is open and bounded, and $\partial U$ is of class $C^1$.
\end{definition}

\begin{theorem}[Representation formula using Green's function]
If $u\in C^2(U)$ solves the Poisson equation
$$
\begin{cases}
-\Delta u=f & \text{in } U,\\
u=g & \text{on } \partial U,
\end{cases}
$$
then
$$
u(x)
=
-\int_{\partial U} g(y)\frac{\partial G}{\partial \nu}(x,y)\,dS(y)
+
\int_U f(y)G(x,y)\,dy.
$$
\end{theorem}

\begin{theorem}[Symmetry of Green's function]
For all $x,y\in U$ with $x\neq y$, we have
$$
G(y,x)=G(x,y).
$$
\end{theorem}

\begin{theorem}[Poisson formula for $\mathbb{R}_+^n$]
Suppose
$
g\in C(\mathbb{R}^{n-1})\cap L^\infty(\mathbb{R}^{n-1}).
$
Let $u:\mathbb{R}_+^n\to \mathbb{R}$ be defined by
$$
u(x)=\int_{\partial \mathbb{R}_+^n} K(x,y)g(y)\,dy
=
\int_{\partial \mathbb{R}_+^n}
\frac{2x_n}{n\alpha(n)}
\frac{1}{|x-y|^n}g(y)\,dy.
$$
Then
\begin{enumerate}
    \item $u\in C^\infty(\mathbb{R}_+^n)\cap L^\infty(\mathbb{R}_+^n)$,
    \item $\Delta u=0$ in $\mathbb{R}_+^n$,
    \item for every $x_0\in \partial \mathbb{R}_+^n$,
    $
    \lim_{x\to x_0}u(x)=g(x_0).
    $
\end{enumerate}
\end{theorem}

\begin{theorem}[Poisson formula for a ball]
Suppose
$
g\in C(\partial B(0,1)).
$
Let $u:B(0,1)\to \mathbb{R}$ be defined by
$$
u(x)=\int_{\partial B(0,1)} K(x,y)g(y)\,dS(y)
=
\int_{\partial B(0,1)}
\frac{1-|x|^2}{n\alpha(n)}
\frac{1}{|x-y|^n}g(y)\,dS(y).
$$
Then
\begin{enumerate}
    \item $u\in C^\infty(B(0,1))$,
    \item $\Delta u=0$ in $B(0,1)$,
    \item for every $x_0\in \partial B(0,1)$,
    $
    \lim_{x\to x_0}u(x)=g(x_0).
    $
\end{enumerate}
\end{theorem}

\section{Mathematical supplement: mollifiers}

\begin{definition}[Standard mollifier]
Define $\eta\in C_c^\infty(\mathbb{R}^n)$ by
$$
\eta(x)=
\begin{cases}
C\exp\!\left(\dfrac{1}{|x|^2-1}\right), & |x|<1,\\[1em]
0, & |x|\ge 1,
\end{cases}
$$
where the constant $C>0$ is selected so that
$$
\int_{\mathbb{R}^n}\eta\,dx=1.
$$
$\eta$ is called the \emph{standard mollifier}.
\end{definition}

\begin{definition}[Mollifier]
For $\varepsilon>0$, set
$$
\eta_\varepsilon(x)=\frac{1}{\varepsilon^n}\eta\!\left(\frac{x}{\varepsilon}\right).
$$
Then the functions $\eta_\varepsilon$ are in $C_c^\infty(\mathbb{R}^n)$ and satisfy
$$
\int_{\mathbb{R}^n}\eta_\varepsilon\,dx=1,
\qquad
\operatorname{supp}(\eta_\varepsilon)\subset B(0,\varepsilon).
$$
\end{definition}

\begin{definition}[Mollification]
If $f:U\to \mathbb{R}$ is locally integrable, define its mollification by
$$
f_\varepsilon=\eta_\varepsilon * f
\qquad \text{in } U_\varepsilon,
$$
that is,
$$
f_\varepsilon(x)
=
\int_U \eta_\varepsilon(x-y)f(y)\,dy
=
\int_{B(0,\varepsilon)} \eta_\varepsilon(y)f(x-y)\,dy.
$$
\end{definition}


\begin{theorem}[Properties of mollifiers]
Here we have the following properties
\begin{enumerate}
    \item $f_\varepsilon\in C^\infty(U_\varepsilon)$.
    \item $f_\varepsilon\to f$ a.e. as $\varepsilon\to 0$.
    \item If $f\in C(U)$, then $f_\varepsilon\to f$ uniformly on compact subsets of $U$.
    \item If $1\le p<\infty$ and $f\in L^p_{\mathrm{loc}}(U)$, then
    $
    f_\varepsilon\to f \text{ in } L^p_{\mathrm{loc}}(U).
    $
\end{enumerate}
\end{theorem}

\section{Regularity and further properties}

\subsection*{Back to Laplace equation}

\begin{theorem}[Smoothness]
If $u\in C(U)$ satisfies the mean value property for each $B(x,r)\subset U$, then
$
u\in C^\infty(U).
$
\end{theorem}

\begin{theorem}[Estimates on derivatives]
Assume $u$ is harmonic in $U$.
Then
$$
|D^\alpha u(x_0)|
\le
\frac{C_k}{r^{n+k}}\|u\|_{L^1(B(x_0,r))}
$$
for each $B(x_0,r)\subset U$ and each multi-index $\alpha$ with order $|\alpha|=k$,
where
$$
C_0=\frac{1}{\alpha(n)},
\qquad
C_k=\frac{(2^n n)^k}{\alpha(n)}.
$$
\end{theorem}

\begin{theorem}[Liouville's theorem]
Suppose $u:\mathbb{R}^n\to \mathbb{R}$ is harmonic and bounded.
Then $u$ is a constant.
\end{theorem}

\begin{theorem}[Representation formula]
For $n\geq 3$, let
$
f\in C_c^2(\mathbb{R}^n).
$
Then any bounded solution of
$
-\Delta u=f \qquad \text{in } \mathbb{R}^3
$
has the form
$$
u(x)=\int_{\mathbb{R}^n}\Phi(x-y)f(y)\,dy + C
$$
for some constant $C$.
\end{theorem}

\begin{remark}
The argument fails for $n=2$, since $\log|x|$ is not bounded at $\infty$.
\end{remark}

\begin{theorem}[Analyticity]
Assume $u$ is harmonic in $U$.
Then $u$ is analytic in $U$.
\end{theorem}